\subsection{内射包络}

定义4. --- 设M为一个A-模。M的一个内射包络是满足以下性质的一对(I, i)，其中I是一个内射A-模，i : M → I是一个同态：

\begin{enumerate}
\def\labelenumi{(\Alph{enumi})}
\setcounter{enumi}{4}
\tightlist
\item
  对于I的一个子模P，P为零当且仅当 \(i^{-1}(\mathbf{P})\) 为零。注意(E)蕴含 \(i\) 是单射的。人们常将M与子模 \(i(\mathbf{M})\) 等同，并称I为M的一个内射包络。
\end{enumerate}

例1. 设A是一个整环，M是无挠的。设K为A的分式域， \(i: \mathbf{M} \to \mathbf{K} \otimes_{\mathbf{A}} \mathbf{M}\) 为典范同态。则 \((\mathbf{K} \otimes_{\mathbf{A}} \mathbf{M}, i)\) 是M的一个内射包络(II, p.～116, 命题26 和 X, p.～17, 推论2)。

定理2. --- 设M为一个A-模。

\begin{enumerate}
\def\labelenumi{\alph{enumi})}
\item
  M具有内射包络。
\item
  如果 (I, i) 和 (J, j) 是 M 的两个内射包络，则存在一个同构 \(f: I \to J\) 使得 \(f \circ i = j\) 。
\end{enumerate}

注意同态 \(f\) （其在 \(b\) 中的存在性被断言）一般不是唯一确定的。

\begin{enumerate}
\def\labelenumi{\alph{enumi})}
\tightlist
\item
  可以假设 M 是内射 A-模 E 的子模（X，第～19 页，推论 3）。考虑按包含关系排序的集合 \(\mathcal{F}\) ，它由 E 中包含 M 且使得典范单射 \(i: \mathrm{M} \to \mathrm{I}\) 满足性质 (E) 的子模 I 组成。由于 \(\mathcal{F}\) 是归纳的，它有极大元（E，III，第～20 页）；设 I 是 \(\mathcal{F}\) 的一个极大元。只需证明 I 是 E 的一个直和项子模。设 N 是 E 的子模，使得 \(\mathrm{N} \cap \mathrm{I} = 0\) ，且对此性质为极大（这样的 N 的存在性见同上所引）。复合同态
\end{enumerate}

\[
\mathrm{I} \xrightarrow {u} \mathrm{E} \xrightarrow {v} \mathrm{E} / \mathrm{N},
\]

其中 \(u\) 和 \(v\) 是典范同态，是单射；由于 \(E\) 是内射的，因此存在一个同态 \(w: E/N \to E\) 使得 \(w \circ v \circ u = u\) ，即 \(w \circ v(x) = x\) （对 \(x \in I\) ）。令 \(I' = \text{Im}(w) = \text{Im}(w \circ v)\) 并设 \(i': M \to I'\) 是典范单射。则 \(I \subset I'\) ；为完成证明，只需证明 \(i'\) 满足条件(E)：这蕴含 \(I = I'\) （I 的极大性），从而 \(w \circ v\) 是 \(E\) 到 \(I\) 上的投影。

设 P 为 I' 的一个子模，使得 \(P \cap M = 0\) 。我们有 \(P = w \circ v(Q)\) ，其中 Q 是 E 的一个包含 N 的子模；此外

\[
\mathbf {Q} \cap \mathbf {M} = w \circ v (\mathbf {Q} \cap \mathbf {M}) \subset \mathbf {P} \cap \mathbf {M} = 0,
\]

因此，由于i：M→I具有性质(E)，Q∩I=0。由N的极大性，这意味着Q=N，即v(Q)=0，从而P=0，正如所要证明的。

\begin{enumerate}
\def\labelenumi{\alph{enumi})}
\setcounter{enumi}{1}
\tightlist
\item
  设 (I, i) 和 (J, j) 是 M 的两个内射包络。由于 J 是内射的，存在一个同态 \(f: I \to J\) 使得 \(f \circ i = j\) ；我们有
\end{enumerate}

\[
i ^ {- 1} (\operatorname{Ker} f) = \operatorname{Ker} j = 0,
\]

因此 Ker f = 0，且 f 是单射。于是 \(f(I)\) 是J的一个内射子模，因此是直和项；由于j满足(E)，这意味着 \(f(I) = J\) 且 f 是双射。

注记. --- 1) 设 (I, i) 是 M 的一个内射包络，并设 j : M → J 是从 M 到内射 A-模 J 的一个单同态。由上述证明可知，存在一个单同态 f : I → J，使得 f ∘ i = j。

\begin{enumerate}
\def\labelenumi{\arabic{enumi})}
\setcounter{enumi}{1}
\tightlist
\item
  设 (I, i) 是 M 的内射包络。将 M 与子模 \(i(\mathbf{M})\) 等同起来；对于I的每一个子模N，存在I的一个内射子模，它是N的内射包络（将注记1应用于N）。反之，I的每一个内射子模J都是子模 \(J \cap M\) 的内射包络。
\end{enumerate}

命题14. --- 设I为非零内射A-模。以下条件等价：

\begin{enumerate}
\def\labelenumi{(\roman{enumi})}
\item
  I是不可分解的（VII，§4，第8号，定义3）；
\item
  0不是I的两个非零子模的交集；
\item
  I 是它所有非零子模的内射包络；
\item
  该环 \(\operatorname{End}_{\mathbf{A}}\) (I) 是局部的（VIII，§ 1，第 4 号，定义 4）。
\item
  \(\Rightarrow\) (iii)：设 M 是 I 的一个非零子模。根据上述注记2，存在 I 的一个子模 I'，它是 M 的内射包络。由于 I' 非零且是 I 的一个直和项，若 I 是不可分解的，则我们有 I = I'。
\item
  \(\Rightarrow\) (ii)：设 E 和 F 是 I 的两个子模，使得 \(\mathbf{E} \cap \mathbf{F} = 0\) 。如果 \(\mathbf{E} \neq 0\) ，则由 (iii) 可知 I 是 E 的内射包络，所以“ \(\mathbf{E} \cap \mathbf{F} = 0\) ”蕴含“ \(\mathbf{F} = 0\) ”.
\item
  \(\Rightarrow\) (i)：这是平凡的。
\item
  \(\Rightarrow\) (i)：这由VIII，§1，第6号，命题13得出。
\item
  \(\Rightarrow\) (iv)：假设I是不可分解的。首先注意，I的每个单射自同态 \(f\) 都是双射（因为此时 \(f(\mathbf{I})\) 是I的一个非零直和项）。此外，每个这样的自同态 \(f\) ——它到 I 的非零子模 E 上的限制是单射的——本身是单射的（事实上，由 (i) \(\Rightarrow\) (iii)，I 是 E 的内射包络，因此“E ∩ Ker \(f = 0\) ”蕴含“Ker \(f = 0\) ”）。鉴于这一点，设 \(f\) 为 \(\operatorname{End}_{\mathbf{A}}(\mathbf{M})\) 的一个不可逆元素；由第八卷第1节第4号命题9可知，只需证明 \(1 - f\) 是可逆的。由于 \(f\) 不是单射的，我们有 \(\operatorname{Ker} f \neq 0\) ；而 \(1 - f\) 到 \(\operatorname{Ker} f\) 上的限制是单射的，故 \(1 - f\) 是单射的，因此是双射的。
\end{enumerate}

推论1. --- 关系“I是不可分解的内射A-模”是集合化的。

事实上，由(iii)可知，每个不可分解的内射A-模都是某个单生成A-模的内射包络。

推论2. --- 设M是一个A-模，并设I是M的内射包络。要使I不可分解，当且仅当0不是M的两个非零子模的交。

该条件由命题 14 的 (i) \(\Rightarrow\) (ii) 可知是必要的。反之，若 I 是两个非零子模的直和 \(\mathbf{I}_1\) 和 \(\mathbf{I}_2\) ，我们将有：

\[
\mathrm{I} _ {1} \cap \mathrm{M} \neq 0, \quad \mathrm{I} _ {2} \cap \mathrm{M} \neq 0 \quad \text { and } \quad (\mathrm{I} _ {1} \cap \mathrm{M}) \cap (\mathrm{I} _ {2} \cap \mathrm{M}) = 0.
\]

例2.——若A是交换且诺特的，则不可分解的内射A-模恰好是模A/p的内射包络，其中p是素理想（X，第171页，习题27）。

